% Generated draft: requires mathematical and source review.
\documentclass[11pt]{article}
\usepackage[margin=25mm]{geometry}
\usepackage{fontspec}
\setmainfont{Latin Modern Roman}
\usepackage{amsmath,amssymb,mathtools,unicode-math}
\setmathfont{Latin Modern Math}
\usepackage{enumitem,needspace,xurl,hyperref}
\setlength{\parindent}{0pt}
\setlength{\parskip}{5pt}
\setlength{\emergencystretch}{3em}
\newcommand{\N}{\mathbb{N}}
\newcommand{\Z}{\mathbb{Z}}
\newcommand{\Q}{\mathbb{Q}}
\newcommand{\R}{\mathbb{R}}
\newcommand{\C}{\mathbb{C}}
\newcommand{\F}{\mathbb{F}}
\newcommand{\A}{\mathbb{A}}
\newcommand{\PP}{\mathbb{P}}
\newcommand{\E}{\mathbb{E}}
\newcommand{\eps}{\varepsilon}
\newcommand{\dd}{\,\mathrm{d}}
\begin{document}
\section*{UnsolvedMath ID: 30005527}
\textbf{Source problem code:} catalogue entry
\subsection{Definitions and mathematical statement}
The supplied record does not specify the ambient group \(G\), the space \(X\), their relationship, or the geometric and rank hypotheses. Consequently, it does not determine a unique fully quantified problem. The following formulation is conditional on the conventional analytic setting described below; these conventions require confirmation against the source context.

Suppose \(G\) is a locally compact unimodular group and \(\Gamma<G\) is discrete. A lattice is such a subgroup for which the invariant quotient measure satisfies \(\operatorname{vol}(\Gamma\backslash G)<\infty\). Thus, under these conventions, a nonlattice satisfies \(\operatorname{vol}(\Gamma\backslash G)=\infty\). Suppose also that \(X\) is a Riemannian \(G\)-space for which the quotient \(\Gamma\backslash X\), its measure, and a nonnegative self-adjoint Laplace operator \(\Delta_{\Gamma\backslash X}\) on \(L^2(\Gamma\backslash X)\) are defined. The intended geometry, treatment of quotient singularities, and operator domain are not supplied.

For an ambient pair \((G,X)\) in the intended, unspecified class, the first question asks whether a nonlattice \(\Gamma<G\) can admit a nonzero square-integrable Laplace eigenfunction, namely whether there can exist
\[
0\ne f\in\operatorname{Dom}(\Delta_{\Gamma\backslash X}),\qquad \lambda\in[0,\infty),\qquad
\Delta_{\Gamma\backslash X}f=\lambda f.
\]
The nonnegative eigenvalue convention here depends on the stated choice of sign for the Laplacian.

For the second question, let \(\rho_\Gamma\) denote the right regular unitary representation on \(L^2(\Gamma\backslash G)\), given by
\[
(\rho_\Gamma(g)F)(\Gamma h)=F(\Gamma h g).
\]
Let \(\lambda_G\) denote the left regular representation on \(L^2(G)\). In the usual weak-containment sense, \(\rho_\Gamma\) is tempered if \(\rho_\Gamma\prec\lambda_G\). Equivalently, for every compactly supported continuous function \(\varphi\) on \(G\),
\[
\left\|\int_G\varphi(g)\rho_\Gamma(g)\,\dd g\right\|
\le
\left\|\int_G\varphi(g)\lambda_G(g)\,\dd g\right\|.
\]
The second question asks whether there exists a nonlattice \(\Gamma<G\) for which \(\rho_\Gamma\not\prec\lambda_G\). The intended restrictions on \(G\) and \(\Gamma\), and whether \(G\) is fixed or varies, remain unresolved. These questions concern the full displayed \(L^2\) spaces; the supplied statement specifies no removal of invariant vectors or restriction to a spectral subspace.
\subsection{Short English statement}
For the intended group \(G\) and geometric space \(X\), can a discrete subgroup of infinite covolume yield a nonzero \(L^2\) Laplace eigenfunction on \(\Gamma\backslash X\)? Can such a subgroup yield a non-tempered right regular representation on \(L^2(\Gamma\backslash G)\)? The ambient hypotheses are missing, so this is a conditional statement-only draft requiring mathematical review.
\subsection{Sources}
\begin{itemize}
\item Supplied local UnsolvedMath registry record, permanent ID \(30005527\), title: \emph{Discrete Spectrum for Nonlattice Quotients}; exact source problem code: \texttt{catalogue entry}.
\item \emph{Group Actions and Harmonic Analysis in Number Theory} (2023), Oberwolfach Reports, pp.\ 1195--1262, DOI \texttt{10.4171/owr/2023/22}: \url{https://doi.org/10.4171/owr/2023/22}. The supplied original excerpt identifies p.\ 1247. The record describes the report as documentation of a workshop and its open-problem discussions. No individual proposer is supplied.
\item Supplied literature reference labeled ``bottom-spectrum theorem'': \url{https://arxiv.org/abs/2304.14565}.
\item Supplied literature reference labeled ``temperedness theorem'': \url{https://arxiv.org/abs/2202.06203}.
\item Literature assessment dated exactly \texttt{2026-08-22}: this is supplied dated curation, not independently verified. It labels the problem partially solved and reports a negative answer to the eigenfunction existence question for nonlattices in higher-rank settings through a bottom-spectrum theorem, together with temperedness for important Anosov classes. It further reports that the full nonlattice classification and possible non-tempered cases remain open. The supplied background explicitly requires independent verification of these time-sensitive claims.
\item The registry labels the statement \texttt{corrected\_verified} and says it was reconstructed and checked against the official report. These are supplied provenance assertions, not verification performed for this draft.
\end{itemize}
\end{document}
